\section{Instrument gate and planned scientific studies}
\label{sec:experiments}

\emph{Status.} This section is the retained historical program description
appended for context; none of the studies below has executed, and the completed
claims of this report are the measurement panels M1--M3 of
\cref{sec:main-studies}.  E3 is deferred by the author's route decision, the
empirical part of E2 is deferred, and the effect-sensitive component of E0-M
remains conditional.  The design has one instrument gate, E0, and five
scientific studies, E1--E5; the disjoint feasibility pilot is a non-inferential
go/no-go stage.  In the
claim ledger E0 remains \emph{not run} until a qualifying completed manifest,
scientific evidence is \emph{none}, and claims are \emph{planned}; smoke,
profiling, and incomplete execution cannot change those statuses.
\Cref{app:analysis} gives the estimators and audits.  The frozen F2-FEHPO report
registers $n_{\mathrm{roots,E3}}=500$, with at most $750$ roots (task templates are
the independent unit); that
count is contested by a corrected simulation and by the adopted verifier
withholding, and an unsigned 150-template, 750-root successor is proposed
(\cref{app:e3-estimator}).  The additional reported
diagnostics of \cref{tab:diagnostics} --- monitor quality, efficiency, learning
behaviour, and generalization axes --- are logged without confirmatory weight.

\paragraph{Two design results before any run.}
The design is already quantitative in a way most pre-empirical papers are not.
Two findings were obtained by simulation before any agent was run, and they
shape the instrument.

\begin{enumerate}[label=(\arabic*),leftmargin=2.2em]
\item \emph{The obvious risk-ranking estimator is badly biased.}  A plug-in
  per-checkpoint effect estimate used for both selection and evaluation returns
  a large positive regret under a null in which the risk score orders checkpoints
  exactly as the verification advantage does: the study would report an
  actionability gap that does not exist.  \Cref{app:e3-estimator} gives the
  split-selection, independent-evaluation estimand that fixes it.
\item \emph{Risk-targeted placement is powerful but changes the tested
  population.}  A semantic fault injected uniformly over eligible checkpoints was
  effectively undetectable in the initial simulation, while lower-risk-tercile
  placement reached $0.98$ at $n=800$, $R=24$.  That placement uses the score
  whose ordering is under test, so it shifts the sampling distribution and can
  restrict score variation.  This does not by itself make a held-out comparison
  circular; reusing outcomes to choose placement and to evaluate would be a
  separate selection problem.  It is retained only as the historical
  F2-ADV adversarial-stress calculation.  The primary F2-FEHPO condition uses risk-blind eligibility and
  randomized matched-versus-mismatched assignment; its power report is
  hash-bound.
\end{enumerate}

Both are calculations, not measurements: they were obtained before any agent was
run.

\paragraph{E0 --- Instrument validation gate.}
\emph{Not a scientific experiment about agents.}  On Tier~0 and then on Tier~1
of \cref{tab:tiers} --- Tiers~2--4 only if they are later registered --- restore
one eligible checkpoint and fork the feasible regimes
into isolated branches; report every layer of \cref{tab:restoration}, and test
cross-branch writes, shared external state, evaluator integrity, feasible-mask
reproducibility, budget restoration, and the cache-independence hypothesis of
\cref{sec:cache} across branch execution positions.  Gate threshold (G1) of
\cref{sec:stages}.  E0 is a gate with no scientific claim about agents attached,
and \cref{thm:safe} depends on it: if branches do not begin from an identical
restored state, no causal estimand here is identified; a failed gate is a
finding about the environment's readiness.

\paragraph{E0-M --- Replay fidelity as a reported measurement.}
\emph{Reported whichever way the gate falls; not a claim about agent
intervention value.}  The runs that produce the E0 verdict also yield per-layer
restoration fidelity over \cref{tab:restoration}, the sibling-leakage rate, and
branch-position and cache effect sizes: these are properties of the isolated
siblings the gate already schedules, and are reported as measurements of the
substrate.  The registered headline quantity, the induced \emph{contrast bias},
is measured on a \emph{separately reserved} pool: siblings forked from a
deliberately unisolated state are compared against correctly isolated siblings,
so the reported number is the regime contrast a study would recover from a state
it only believed it had restored.  It cannot come from the E0 runs, because the
adopted instrument pool prohibits any effect estimate and a deliberately
unisolated condition belongs to its own registered pool and namespace; the
reservation is template indices 420--479, and its size is adopted after the
pilot informs feasibility and precision.  \emph{Why it is separated from E0:} the gate
answers whether this substrate may carry causal claims, whereas E0-M answers how
much a replay-based design is wrong by when it is wrong --- a quantity the
branching and counterfactual-replay work of \cref{sec:related} assumes rather
than reports.  It licenses no agent-effect claim and cannot substitute for E1--E5.

\paragraph{E1 --- Representation adequacy.}  \emph{Tied to
\cref{thm:refinement}.}  Does the registered $Z=f(H)$ discard decision-relevant
information?  Estimate $V^{*Z'}_j-V^{*Z}_j$ on held-out branches for a richer
registered $Z'$, which by \eqref{eq:refinement-bound} lower-bounds the
representation loss.  Comparators: baseline risk alone; elapsed steps; risk
with steps; full history.  Beside the value difference, report
representation-dependent policy value and \emph{action-choice disagreement} ---
the fraction of eligible checkpoints at which the $Z$- and $Z'$-optimal rules
select different regimes --- at matched availability and matched expected cost,
so that a difference in decisions is separated from a difference in value.  Any
$Z'$, diagnostic, or selection criterion is declared before its evaluation data
are read.  E0 and E1 answer independent questions: a sufficient representation
does not repair a corrupted restore, and an exact restore does not make a
compressed observation sufficient.  \emph{Prediction:} $Z$ strictly beats
risk-with-elapsed-steps.  \emph{Falsification:} if not, $Z$ is not earning its
complexity and representation design is revisited before any policy machinery.
A null result bounds the value of the richer registered $Z'$ on this
population; it is not evidence that $Z$ is Markov or conditional-mean
sufficient.

\paragraph{E2 --- The actionability gap.}  \emph{Deferred characterization;
only the availability identity is retained as explanatory algebra
(\cref{sec:main-measurements}); the empirical study is deferred.}  It would
estimate the joint distribution of baseline
failure risk and actionable value \eqref{eq:actionable}, stratified by fault
family, propagation depth, remaining budget, availability, and verifier
reliability, and report the mass in the three regimes of \cref{sec:policy}
together with the $J=0$ harm fraction \eqref{eq:harm-fraction}.
\emph{Prediction:} all three regimes carry non-negligible mass and regime~(iii)
is non-empty in the no-fault stratum.  \emph{Primary prediction:} the gap
appears under F2-FEHPO on the approved task and tier population.  F1-ORD is a
historical mechanically observable contrast and F2-ADV a historical
adversarial-stress calculation; neither is executed in the current protocol or
counts as replication of the positive mechanism.  A one-tier gap is reported as such.
\emph{Falsification:} if actionable value is monotone in baseline
risk, the gap is empty and the actionable-value framing collapses to risk ranking.

\paragraph{E3 --- Ordinal sufficiency.}  \emph{Deferred; not the current
primary claim; tied to \cref{thm:ordinal}.}  The design would test the null
\eqref{eq:null} by comparing budget-matched
capacity-$q$ allocation under a cross-fitted effect score against the risk
score, on the joint $g_0$/$g_v$ domain, across a declared capacity grid.  The
estimator, and why the naive plug-in is invalid, are in \cref{app:e3-estimator}.
The comparison is between the ordering induced by baseline failure risk
$1-V^\rho$ and the ordering induced by the advantage $A^\rho(\cdot,g_v)$, at a
fixed \emph{count} of verifications per capacity.  If expected verification cost
is heterogeneous, the corresponding resource-constrained problem is
$\max_u\E[u(W)a_v(W)]$ subject to $\E[u(W)c_v(W)]\le B$ over $0\le u\le1$, whose
relaxed solution thresholds $a_v-\lambda c_v$; that quantity is \emph{proposed}
in \cref{app:analysis} as an optional secondary analysis --- it has no
registration event behind it --- and is never an E3 endpoint.
\emph{Prediction:} the null is rejected at some capacity under F2-FEHPO.
\emph{Why it was primary:} it is the sharp null and kill criterion that would
control the insufficiency framing.  It needs no recovery regime and therefore
has a distinct feasibility and power calculation from the E4 safety study.  Its approved
root count $n_{\mathrm{roots,E3}}=500$ (maximum authorized $750$) comes from the
prospective F2-FEHPO report, not from the E4
certification ceiling \eqref{eq:sample-ceiling}.
\emph{Falsification (kill criterion):} if ordinal insufficiency is not
established under the approved F2-FEHPO design, the insufficiency framing is
withdrawn and simultaneous intervals and power are reported.  Non-rejection
does not establish ordinal sufficiency; historical contrasts cannot rescue the claim.

\paragraph{E4 --- Certified safe improvement at matched cost.}  \emph{Safety claim; tied to
\cref{thm:safe} and \cref{prop:certification}.}  Does certification add safety
beyond constrained training, and what does that governance cost?  E4 reports
the pre-registered decomposition of the policy-learning arms of \cref{sec:policy},
\begin{equation}
 \begin{aligned}
  \Delta_{\mathrm{learning}}&=V(d_{\mathrm{U}})-V(d^\star_{\mathrm{fixed}}),\\
  \Delta_{\mathrm{constraint}}&=V(d_{\mathrm{C}})-V(d_{\mathrm{U}}),\\
  \Delta_{\mathrm{certification}}&=V(d_{\mathrm{C,cert}})-V(d_{\mathrm{C}}),
 \end{aligned}
 \label{eq:e4-decomposition}
\end{equation}
where $d_{\mathrm{U}}$, $d_{\mathrm{C}}$, and $d_{\mathrm{C,cert}}$ are the
unconstrained, constrained, and certified constrained policies.  Across the prevalence grid,
with the no-fault harm fraction \eqref{eq:harm-fraction} and the fallback rate
of \eqref{eq:gated-rule} as primary endpoints, the last two differences are
expected to be negative in mean completion and positive in safety; that is the
point, because the decomposition prices governance rather than assuming it is
free.  Completion, no-fault harm, invalidity, every constrained cost component,
and fallback are reported at comparable resource envelopes, and
$\Delta_{\mathrm{certification}}$ is the empirical contribution of uncertainty
handling and of nothing else.  Population constraint satisfaction is not a hard
per-branch ceiling: \eqref{eq:cost-constraint} bounds an expectation, and the
per-regime hard ceilings of \cref{sec:constraints} are operational rails whose
target binding rate is zero.  \emph{Falsification:} if $\Delta_{\mathrm{certification}}\approx0$ on
the no-fault harm fraction, the certification framework adds nothing beyond
constrained training and the paper reports that; if the optimistic and gated
controllers perform identically, the lower confidence bound is doing no work and
\cref{thm:safe} is decoration.

\paragraph{E5 --- Verifier independence under attack.}  \emph{Robustness study.}
The regime $g_v$ invokes
a registered independent verifier, and independence is an assertion about an
engineered system: corrupt the verifier channel at controlled rates and
separately correlate it with the controller's selection signal, measuring
misestimation of $\mu_{g_v}$, deformation of the deviation set $\mathcal N$, and
whether the fallback rate rises.  This perturbs the verification channel while
terminal adjudication stays independent; a benchmark that adds task-irrelevant
context to the agent's input \citep{asibench2026} perturbs a different
channel, and the two motivate distinct robustness checks.  \emph{Prediction:} across the declared
corruption grid the gated controller's harm fraction rises more slowly than the
optimistic controller's at matched cost, and that fallback rises under
undetected corruption.  Both are predictions under test: an undetected
corruption does not automatically produce more fallback, and a
clean-calibration interval is not valid after arbitrary verifier shift unless
the registered uncertainty set covers that shift.  Where the design permits,
independent corruption is separated from corruption correlated with the
controller's selection signal, and the controller-visible verifier output is
kept distinct from terminal adjudication throughout.
\emph{Falsification:} if the two are
indistinguishable on that grid, pessimism buys nothing once the independence
assertion fails, and the gate is reported as an unnecessary cost rather than
defended.  This is the argument for pessimism that does not depend on
\cref{thm:safe} holding exactly.  \emph{What E5 is not:} corrupting the channel
does not demonstrate reward hacking.  That would require additionally
optimising a controller \emph{against} the flawed verifier score and evaluating
it against an independent target, which is the proxy-versus-true-performance
divergence studied by \citet{pan2022} and is not part of this protocol.

\paragraph{E3-design feasibility pilot.}  Between E0 and the confirmatory study, a pilot
estimates eligible-checkpoint yield, restoration and instrumentation
reliability, availability per regime, cost distributions, adjudicator agreement,
between- and within-checkpoint variance, and whether recovery semantics remain
coherent under each family.  Pilot data are excluded from confirmatory
inference; the pilot reports whether F2-FEHPO can support the approved E3 design
at the projected cost.  Its roots, tasks, instances, and seeds are excluded from
confirmatory evaluation, and it cannot satisfy or kill E3 or tune the fault
definition and capacity grid to the observed effect.
